\documentclass[11pt]{article}

\usepackage[margin=1in]{geometry}
\usepackage{booktabs}
\usepackage{array}
\usepackage{tabularx}
\usepackage{longtable}
\usepackage{graphicx}
\usepackage{subcaption}
\usepackage{float}
\usepackage{xcolor}
\usepackage{hyperref}

\hypersetup{
  colorlinks=true,
  linkcolor=blue!60!black,
  citecolor=blue!60!black,
  urlcolor=blue!60!black
}

\title{\textbf{Market Making with Deep Reinforcement Learning from Limit Order Books}\\
\large Reproduction, Data Adaptation, and Robustness Extensions}
\author{Jiaan Han \and Xiangpeng Huang \and Tianle Liu}
\date{}

\newcommand{\szroot}{Our-Code/experiment_results/sz10_full_experiment}
\newcommand{\lobroot}{Our-Code/experiment_results/lobster3_full_experiment}
\newcommand{\codepath}[1]{\path{#1}}

\begin{document}
\maketitle

\begin{abstract}
This project studies a deep reinforcement learning approach to market making from
limit order book data. We run experiments on two market data sources: LOBSTER
NASDAQ order book data for AAPL, AMZN, and MSFT, and Shenzhen Stock Exchange
Level-2 data for 10 selected Chinese stocks. We begin from the paper
\emph{Market Making with Deep Reinforcement Learning from Limit Order Books},
which uses an Attn-LOB representation model before training a reinforcement
learning market maker. We reimplemented the pipeline in PyTorch, added reusable
training and evaluation interfaces, reconstructed a Shenzhen Level-2 data set
from order and trade events, and compared reinforcement learning agents with
fixed-level, random, and Avellaneda-Stoikov baselines. We also added robustness
experiments that separate observation latency, order-arrival latency, order
loss, and model inference latency. Our strongest classification result is a
joint 10-stock Attn-LOB encoder with test macro F1 of 0.701 on Shenzhen data. In
trading backtests, a simple fixed-level baseline is strongest on the Shenzhen
sample, while the joint Dueling-DQN agent is strongest on the three-stock
LOBSTER sample. These results suggest that the learned representation is useful,
but that profitable learned market making remains highly sensitive to the data
source, execution model, latency assumptions, and inventory control.
\end{abstract}

\section{Introduction}

Market making is the problem of continuously posting bid and ask quotes, earning
spread revenue while managing inventory and adverse selection risk. In this
project, we evaluate this problem on both LOBSTER NASDAQ data and Shenzhen Stock
Exchange Level-2 data, so that the same modeling pipeline can be tested on the
original-style U.S. order book format and on a reconstructed Chinese limit order
book data set. The task is well suited to reinforcement learning (RL) because
actions affect future inventory, realized PnL, and risk exposure. It is also
difficult: the agent observes a high-dimensional, rapidly changing limit order
book (LOB), while executions depend on queue position, latency, liquidity, and
short-horizon price movement.

The original paper we studied proposes a two-stage architecture. First, an
Attn-LOB neural network is trained on limit order book states to predict
short-horizon mid-price movement. Second, the learned encoder is used inside a
market-making reinforcement learning agent. Our project reproduces this idea and
extends it in four directions:

\begin{enumerate}
  \item We rewrote the Attn-LOB model and training loop in PyTorch.
  \item We built a full data pipeline for Shenzhen Level-2 order and trade data,
  reconstructing an event-level 10-level order book.
  \item We evaluated PPO and Dueling-DQN against interpretable non-learning
  baselines.
  \item We added robustness tests for observation latency, order latency, order
  loss, and inference latency, plus regime-sensitivity analysis.
\end{enumerate}

All final experiment outputs are stored locally under
\codepath{Our-Code/experiment_results/}. The code interface is documented in
\codepath{Our-Code/README.md}, \codepath{Our-Code/SYSTEM_README.md}, and
\codepath{Our-Code/PROJECT_RUNBOOK.md}.

\section{Original Paper and Reproduction Target}

The original paper combines representation learning and reinforcement learning.
The representation stage uses a convolutional and attention-based LOB model to
predict whether the future mid-price moves down, stays stationary, or moves up.
This supervised task is not the final trading objective, but it encourages the
encoder to extract short-horizon price formation signals from the LOB. The RL
stage then uses that encoder to select market-making quotes.

Our first reproduction target was a one-day LOBSTER AAPL experiment. During the
initial model search, we compared the original Attn-LOB configuration with wider
architectures and different optimization choices. Table~\ref{tab:aapl_tuning}
summarizes the most relevant early pretraining results from our project slides.
The best modified setting used AdamW and weighted cross entropy with embedding,
convolution, and inception dimensions of 96/32/96.

\begin{table}[H]
\centering
\caption{Initial AAPL Attn-LOB tuning results from the reproduction stage.}
\label{tab:aapl_tuning}
\begin{tabular}{llccccc}
\toprule
Setting & Optimizer + loss & Size & Params & Precision & Recall & F1 \\
\midrule
Original & Adam + CE & 64/32/64 & 271k & 0.438 & 0.454 & 0.434 \\
Best architecture-only & Adam + CE & 128/48/128 & 613k & 0.461 & 0.458 & 0.453 \\
Best modified & AdamW + weighted CE & 96/32/96 & 377k & 0.505 & 0.518 & 0.510 \\
\bottomrule
\end{tabular}
\end{table}

We used this best modified Attn-LOB configuration for the larger experiments in
this report.

\section{Data Sets}

\subsection{LOBSTER Sample Data}

The first data source is a LOBSTER sample containing AAPL, AMZN, and MSFT for
June 21, 2012. LOBSTER already provides message files and 10-level order book
snapshots. We convert the files into a common event format used by both
pretraining and RL. The sample is useful for reproducing the original data
format, but it is small: each stock has only one trading day.

\subsection{Shenzhen Level-2 Data}

The second data source is a five-day Shenzhen Level-2 data set covering January
21, 22, 23, 24, and 27, 2025. Unlike LOBSTER, the uploaded data include separate
market snapshot, order, and trade HDF5 files. We therefore built a replay
pipeline:

\begin{enumerate}
  \item read order submissions and trade/cancel records;
  \item maintain active buy and sell order maps by order id;
  \item aggregate active orders into price levels;
  \item emit 10-level LOB states and event features;
  \item validate reconstructed best bid/ask against market snapshots;
  \item compute train/validation/test labels from future mid-price movement.
\end{enumerate}

We selected 10 Shenzhen stocks using only data-quality criteria: completeness
across dates, activity, spread, valid-book fraction, label balance, and
non-stationary signal. The final split is date based: January 21--23 for
training, January 24 for validation, and January 27 for testing. This prevents
future-day leakage. The processed event files contain LOB depth, execution
features, split id, and labels in $\{-1,0,1\}$.

\begin{table}[H]
\centering
\caption{Main experiment data sets.}
\label{tab:datasets}
\begin{tabularx}{\textwidth}{lXccc}
\toprule
Data set & Securities & Dates & Split & Main use \\
\midrule
SZ10 & 10 Shenzhen stocks & 5 days & 3:1:1 by date & Main large-scale experiment \\
LOBSTER3 & AAPL, AMZN, MSFT & 1 day & intraday split & Original-format comparison \\
\bottomrule
\end{tabularx}
\end{table}

\section{Methodology}

\subsection{Attn-LOB Pretraining}

The supervised pretraining target is the three-class future mid-price movement
label. Given a 50-event LOB window, the model predicts down, stationary, or up.
The encoder contains convolutional feature extraction across price/size columns,
inception-style temporal filters, and multi-head attention. We train either:

\begin{itemize}
  \item a \textbf{single-stock encoder}, trained separately for each stock; or
  \item a \textbf{joint encoder}, trained once on all selected stocks and reused
  by per-stock RL agents.
\end{itemize}

Joint pretraining is meant to test whether shared order-book structure transfers
across securities.

\subsection{Reinforcement Learning Environment}

The RL environment is an event-driven market-making simulator. At each decision
step, the agent observes normalized LOB windows, auxiliary event features, and
inventory state. The environment tracks cash, inventory, mark-to-market value,
trading volume, average spread, realized PnL, normalized PnL, PnL per map, and
inventory penalties.

We evaluate two agent families:

\begin{itemize}
  \item \textbf{PPO}, using a continuous action that controls reservation price
  bias and quote spread.
  \item \textbf{Dueling-DQN}, using a discrete action space for quote placement
  and inventory management.
\end{itemize}

The reward combines dampened PnL, trading PnL, and an inventory penalty. All
reported results use zero explicit transaction cost and zero market impact, which
matches our small-order, high-frequency market-making assumption. The code still
contains transaction-cost parameters for future sensitivity checks.

\subsection{Baselines}

We compare the RL agents with non-learning quote policies:

\begin{itemize}
  \item fixed level 1, 2, or 3: quote at the corresponding LOB level;
  \item random top-5-level quoting: sample bid and ask levels independently;
  \item Avellaneda-Stoikov: quote around an inventory-adjusted reservation price.
\end{itemize}

These baselines are important because raw RL PnL is difficult to interpret
without a simple market-making benchmark.

\subsection{Latency, Order Loss, and Regime Tests}

The original paper-style latency is an observation delay: the agent observes
state $t-\ell$ when quoting at event $t$. We extended this into several
execution-friction modes:

\begin{itemize}
  \item observation latency in events;
  \item order-arrival latency in events or milliseconds;
  \item independent bid/ask order loss;
  \item simulated inference latency added to order arrival;
  \item measured model decision time.
\end{itemize}

We also group test episodes into low, mid, and high regimes by realized
volatility, average spread, and trade intensity.

\section{Implementation}

The implementation is organized as a reusable system rather than a collection of
one-off scripts. The main entry point is:

\begin{verbatim}
cd Our-Code
../.venv/bin/python3 run_system.py <command>
\end{verbatim}

The full SZ10 pipeline can be run or resumed with:

\begin{verbatim}
../.venv/bin/python3 run_system.py full --stage all
\end{verbatim}

Individual stages include \codepath{pretrain}, \codepath{rl},
\codepath{baseline}, \codepath{latency}, \codepath{regime},
\codepath{summarize}, and \codepath{plots}. The LOBSTER3 comparison can be
reproduced with:

\begin{verbatim}
../.venv/bin/python3 run_lobster3_full_experiment.py --stage all
\end{verbatim}

The data preparation scripts are grouped under \codepath{Our-Code/data_prep/}.
The API wrapper is in \codepath{Our-Code/system_api.py}.

\section{Results}

\subsection{Pretraining Results}

Table~\ref{tab:pretrain_main} shows the main classification results. Joint
pretraining is especially strong on the Shenzhen data, reaching 0.701 test macro
F1. On LOBSTER3, the joint encoder also improves macro F1 relative to the
single-stock encoders.

\begin{table}[H]
\centering
\caption{Attn-LOB pretraining results.}
\label{tab:pretrain_main}
\begin{tabular}{llrrrr}
\toprule
Data & Scope & Accuracy & Precision & Recall & F1 \\
\midrule
SZ10 & Best single stock (002193.SZ) & 0.725 & 0.655 & 0.668 & 0.658 \\
SZ10 & Joint 10-stock encoder & 0.706 & 0.702 & 0.701 & 0.701 \\
LOBSTER3 & Best single stock (MSFT) & 1.000 & 0.333 & 0.333 & 0.333 \\
LOBSTER3 & Joint 3-stock encoder & 0.658 & 0.485 & 0.692 & 0.497 \\
\bottomrule
\end{tabular}
\end{table}

\begin{figure}[H]
\centering
\includegraphics[width=0.82\textwidth]{\szroot/final_plots/03_pretrain_test_f1.pdf}
\caption{SZ10 pretraining F1 comparison for single-stock and joint encoders.}
\label{fig:sz_pretrain}
\end{figure}

\subsection{Trading Results on SZ10}

Table~\ref{tab:sz_trading} reports test metrics averaged across the 10 selected
Shenzhen stocks. Fixed level 2 is the best method by mean PnL and Sharpe-like
stability, followed by fixed level 3 and joint PPO. This is an important result:
although the joint encoder predicts price movement well, the learned RL policy
does not dominate simple spread-capture baselines on this short sample.

\begin{table}[H]
\centering
\caption{SZ10 aggregate trading results on the test day.}
\label{tab:sz_trading}
\small
\begin{tabular}{lrrrrr}
\toprule
Method & PnL & ND-PnL & PnL-MAP & Avg abs inv. & Volume \\
\midrule
baseline\_fixed\_2 & 13.97 & 406.36 & 0.110 & 236.60 & 9,198 \\
baseline\_fixed\_3 & 6.57 & 103.83 & 0.069 & 91.69 & 2,424 \\
joint\_ppo & 3.50 & 17.63 & 0.081 & 152.72 & 9,715 \\
baseline\_as & 0.54 & 4.61 & 0.027 & 50.79 & 1,412 \\
baseline\_random\_5level & 0.09 & 9.25 & 0.016 & 314.66 & 17,071 \\
single\_ppo & -2.22 & -48.30 & -0.020 & 142.42 & 13,189 \\
joint\_dueling\_dqn & -8.83 & -258.03 & -0.020 & 240.29 & 32,308 \\
single\_dueling\_dqn & -9.85 & -279.50 & -0.028 & 269.65 & 31,059 \\
baseline\_fixed\_1 & -20.62 & -1355.77 & -0.021 & 508.06 & 61,707 \\
\bottomrule
\end{tabular}
\end{table}

\begin{figure}[H]
\centering
\includegraphics[width=0.82\textwidth]{\szroot/final_plots/01_overall_test_pnl.pdf}
\caption{SZ10 overall test PnL comparison.}
\label{fig:sz_pnl}
\end{figure}

\begin{figure}[H]
\centering
\includegraphics[width=0.82\textwidth]{\szroot/final_plots/02_risk_return_inventory.pdf}
\caption{SZ10 risk, return, and inventory diagnostics.}
\label{fig:sz_risk}
\end{figure}

\subsection{Trading Results on LOBSTER3}

Table~\ref{tab:lob_trading} reports the AAPL/AMZN/MSFT comparison. Here the
ranking changes: joint Dueling-DQN has the best mean PnL, while PPO performs
poorly. The difference between SZ10 and LOBSTER3 shows that agent rankings are
data-dependent and that a policy that is viable in one market sample may not
generalize automatically.

\begin{table}[H]
\centering
\caption{LOBSTER3 aggregate trading results on the test split.}
\label{tab:lob_trading}
\small
\begin{tabular}{lrrrrr}
\toprule
Method & PnL & ND-PnL & PnL-MAP & Avg abs inv. & Volume \\
\midrule
joint\_dueling\_dqn & 11.03 & 46.08 & 0.026 & 233.72 & 288,218 \\
baseline\_fixed\_2 & 0.00 & 0.00 & 0.000 & 0.00 & 0 \\
baseline\_fixed\_3 & 0.00 & 0.00 & 0.000 & 0.00 & 0 \\
baseline\_random\_5level & -2.67 & -23.96 & 0.021 & 123.35 & 122,066 \\
baseline\_fixed\_1 & -9.70 & -560.94 & 0.003 & 464.16 & 481,241 \\
baseline\_as & -11.77 & -114.19 & 0.006 & 404.00 & 526,865 \\
single\_dueling\_dqn & -30.60 & -322.56 & -0.061 & 333.30 & 277,048 \\
joint\_ppo & -79.57 & -1098.72 & -0.121 & 406.74 & 844,747 \\
single\_ppo & -99.07 & -1692.37 & -0.149 & 406.73 & 900,857 \\
\bottomrule
\end{tabular}
\end{table}

\begin{figure}[H]
\centering
\includegraphics[width=0.82\textwidth]{\lobroot/final_plots/01_overall_test_pnl.pdf}
\caption{LOBSTER3 overall test PnL comparison.}
\label{fig:lob_pnl}
\end{figure}

\subsection{Single vs Joint Encoder}

The results support joint representation learning more clearly than they support
a universal RL policy improvement. On SZ10, the joint encoder is much stronger
than the median single-stock encoder in classification, and joint PPO is the
best RL agent. On LOBSTER3, the joint encoder also improves classification and
joint Dueling-DQN is the best trader. However, joint Dueling-DQN underperforms on
SZ10, which suggests that the downstream RL objective is still sensitive to
reward shape, action space, and execution details.

\begin{figure}[H]
\centering
\includegraphics[width=0.82\textwidth]{\szroot/final_plots/04_per_stock_pnl_heatmap.pdf}
\caption{Per-stock SZ10 PnL heatmap. Performance varies substantially by stock.}
\label{fig:sz_heatmap}
\end{figure}

\section{Robustness Results}

\subsection{Latency and Order Loss}

Table~\ref{tab:latency} summarizes representative latency scenarios. The SZ10
fixed-level baseline is relatively robust, while joint PPO remains positive in
several observation-latency and order-loss scenarios. Joint Dueling-DQN is more
fragile on SZ10 but remains the best learning agent on LOBSTER3 under several
observation-latency scenarios.

\begin{table}[H]
\centering
\caption{Representative latency and order-loss PnL results.}
\label{tab:latency}
\small
\begin{tabular}{llrrrr}
\toprule
Data & Scenario & fixed\_2 & AS & joint PPO & joint DQN \\
\midrule
SZ10 & paper\_base & 13.97 & 0.54 & 3.50 & -8.83 \\
SZ10 & obs\_latency\_5 & 18.45 & 6.14 & 5.87 & -3.79 \\
SZ10 & obs\_latency\_10 & 2.40 & 6.23 & -0.21 & -23.90 \\
SZ10 & loss\_5pct & 14.00 & 0.76 & 4.42 & -5.90 \\
SZ10 & infer50\_loss5pct & 2.27 & 1.84 & -6.00 & -20.19 \\
\midrule
LOBSTER3 & paper\_base & 0.00 & -11.77 & -79.57 & 11.03 \\
LOBSTER3 & obs\_latency\_5 & 10.20 & -18.63 & -88.33 & 25.37 \\
LOBSTER3 & obs\_latency\_10 & 15.53 & -25.63 & -99.07 & 16.63 \\
LOBSTER3 & loss\_5pct & 0.00 & -15.90 & -75.03 & 14.60 \\
LOBSTER3 & infer50\_loss5pct & 2.60 & -26.33 & -29.73 & -15.53 \\
\bottomrule
\end{tabular}
\end{table}

\begin{figure}[H]
\centering
\includegraphics[width=0.82\textwidth]{\szroot/final_plots/05_latency_pnl_delta.pdf}
\caption{SZ10 PnL change under latency, order-loss, and inference-delay scenarios.}
\label{fig:latency_delta}
\end{figure}

\subsection{Regime Sensitivity}

Regime analysis groups episodes by realized volatility, average spread, and
trade intensity. For SZ10, fixed level 2 performs best in high-volatility,
high-spread, and high-intensity regimes. Joint PPO is positive in high-volatility
and high-spread regimes but weaker in low-volatility episodes. For LOBSTER3,
joint Dueling-DQN performs well in high-volatility and low trade-intensity
episodes, while PPO is consistently negative.

\begin{figure}[H]
\centering
\includegraphics[width=0.82\textwidth]{\szroot/final_plots/07_regime_sensitivity_pnl.pdf}
\caption{SZ10 regime-sensitivity results.}
\label{fig:regime}
\end{figure}

\section{Discussion}

Our main findings are:

\begin{enumerate}
  \item \textbf{Attn-LOB pretraining is useful.} The joint Shenzhen encoder
  achieves strong out-of-sample classification performance, suggesting that
  cross-stock representation learning captures shared microstructure patterns.
  \item \textbf{Better prediction does not automatically imply better market
  making.} Fixed-level quotes can outperform RL because the trading objective
  depends on inventory, fill probability, and spread capture, not only direction
  prediction.
  \item \textbf{Baselines are essential.} Without fixed-level, random, and
  Avellaneda-Stoikov benchmarks, RL PnL is hard to interpret.
  \item \textbf{Execution assumptions change conclusions.} Observation latency,
  order latency, order loss, and inference delay can change method rankings.
  \item \textbf{Market and sample dependence is large.} SZ10 and LOBSTER3 lead
  to different best agents, so robustness across assets and days remains an open
  issue.
\end{enumerate}

There are also limitations. The simulator does not model exact exchange queue
position, partial priority, hidden liquidity, or market impact. We set explicit
transaction costs to zero in the main experiments because the intended strategy
is small-order market making, but real deployments would need exchange-specific
fees, rebates, and latency distributions. The Shenzhen sample has only five
trading dates, so the learned policies should be treated as experimental rather
than production-ready.

\section{Conclusion}

We built a complete reproduction and extension of an Attn-LOB based
market-making framework. The project now includes data reconstruction,
pretraining, RL training, baseline evaluation, latency and order-loss stress
tests, regime sensitivity, and reproducible command-line interfaces. The strongest
classification model is the joint SZ10 Attn-LOB encoder. The strongest SZ10
trading method is the fixed-level 2 baseline, while the strongest LOBSTER3
trading method is joint Dueling-DQN. The difference between these results is
itself informative: market-making RL needs careful evaluation against simple
baselines and realistic execution frictions before its PnL can be interpreted.

\appendix
\section{Reproducibility Commands}

\begin{verbatim}
cd Our-Code

# Rebuild reports and final plots from saved results
../.venv/bin/python3 run_system.py summarize

# Run the full SZ10 pipeline
../.venv/bin/python3 run_system.py full --stage all

# Run latency scenarios
../.venv/bin/python3 run_system.py latency --eval-episodes 10

# Run regime sensitivity
../.venv/bin/python3 run_system.py regime

# Run the LOBSTER3 comparison
../.venv/bin/python3 run_lobster3_full_experiment.py --stage all
\end{verbatim}

\section{Output Files}

The most important output files are:

\begin{itemize}
  \item SZ10 root: \codepath{Our-Code/experiment_results/sz10_full_experiment/}
  \item SZ10 reports: \codepath{reports/pretrain_summary.csv},
  \codepath{reports/comparison_test_aggregate.csv}
  \item SZ10 robustness: \codepath{latency_scenarios/reports/latency_aggregate.csv},
  \codepath{regime_sensitivity/reports/regime_sensitivity_summary.csv}
  \item SZ10 figures: \codepath{final_plots/}
  \item LOBSTER3 root: \codepath{Our-Code/experiment_results/lobster3_full_experiment/}
  \item LOBSTER3 reports and figures:
  \codepath{reports/comparison_test_aggregate.csv}, \codepath{final_plots/}
\end{itemize}

\begin{thebibliography}{9}

\bibitem{original}
Original project paper: \emph{Market Making with Deep Reinforcement Learning
from Limit Order Books}.

\bibitem{avellaneda}
Avellaneda, M. and Stoikov, S. High-frequency trading in a limit order book.
\emph{Quantitative Finance}, 2008.

\bibitem{ppo}
Schulman, J., Wolski, F., Dhariwal, P., Radford, A., and Klimov, O.
Proximal Policy Optimization Algorithms. 2017.

\bibitem{dueling}
Wang, Z., Schaul, T., Hessel, M., van Hasselt, H., Lanctot, M., and de Freitas,
N. Dueling Network Architectures for Deep Reinforcement Learning. 2016.

\end{thebibliography}

\end{document}
